\documentclass[11pt,a4paper]{article}
\usepackage[utf-8]{inputenc}
\usepackage{amsmath}
\usepackage{amssymb}
\usepackage{geometry}
\usepackage{hyperref}
\usepackage{graphicx}
\usepackage{booktabs}
\usepackage{array}

\geometry{margin=1in}
\setlength{\parindent}{0pt}
\setlength{\parskip}{0.5em}

\title{Integrated Terrain-Aware Navigation and Locomotion for Quadruped Robots using Hierarchical MPC}
\author{}
\date{}

\begin{document}

\maketitle

\section{SYSTEM ARCHITECTURE}

We present a tightly-integrated locomotion and navigation framework for the Unitree Go2 quadruped robot operating on unstructured terrain. The key insight is that a single terrain representation---a dual-layer heightmap constructed from onboard LiDAR---simultaneously informs decisions at every level of the control hierarchy: global path planning, local trajectory optimization, contact force generation, foothold selection, and swing leg control.

The system operates as a cascaded loop at four rates:

\begin{center}
\begin{tabular}{|c|c|l|}
\hline
\textbf{Layer} & \textbf{Rate} & \textbf{Function} \\
\hline
Physics simulation & 1000 Hz & MuJoCo rigid-body dynamics \\
Leg controller & 200 Hz & Cartesian impedance + force tracking \\
Centroidal MPC & 30--50 Hz & Convex QP for optimal contact forces \\
MPPI planner & 15--25 Hz & Sampling-based velocity command generation \\
Perception + A* & $\sim$6 Hz & LiDAR scan, heightmap update, global replan \\
\hline
\end{tabular}
\end{center}

The terrain representation is consumed by five distinct subsystems:
\begin{enumerate}
\item Friction cone rotation in the MPC
\item COM height adaptation in the trajectory generator
\item Foothold quality scoring in the gait planner
\item Steppability and obstacle costs in the MPPI local planner
\item Edge traversability in the A* global planner
\end{enumerate}

This coupling ensures that the robot's navigation decisions are grounded in the same physical constraints that govern its locomotion.

\section{TERRAIN PERCEPTION AND REPRESENTATION}

\subsection{3D LiDAR Simulation}

We simulate a 3D LiDAR sensor mounted at the robot's trunk using MuJoCo ray-casting. The sensor fires $N_{az} \times N_{el} = 90 \times 15 = 1350$ rays per scan, spanning azimuth $[-\pi, \pi)$ and elevation $[-30°, 15°]$ with a maximum range of 6 m. Each ray direction in the body frame is

\begin{equation}
\mathbf{d}_{\text{body}} = \begin{bmatrix} \cos\phi\cos\theta \\ \cos\phi\sin\theta \\ \sin\phi \end{bmatrix}
\end{equation}

where $\theta$ is azimuth and $\phi$ is elevation. Rays are rotated to world frame via the trunk orientation $\mathbf{R}_{\text{wb}}$, and hits on the robot's own body are discarded. Points with radial distance $< 0.45$ m from the base or with height outside $[-0.5, 2.0]$ m are filtered.

\subsection{Dual-Layer Global Heightmap}

We maintain a global heightmap $\mathcal{H}$ on a 2D grid of resolution $\delta = 0.05$ m covering a $12 \times 12$ m region. Each cell $(i,j)$ stores two height values updated via exponential moving average (EMA):

\textbf{Ground layer} (walkable surface estimate):
\begin{equation}
h^{\text{ground}}_{i,j} \leftarrow (1 - \alpha_g)\, h^{\text{ground}}_{i,j} + \alpha_g\, Q_{0.2}(\{z_k\}_{k \in \text{cell}})
\end{equation}

where $Q_{0.2}$ denotes the 20th percentile of height samples within the cell, and $\alpha_g = 0.25$.

\textbf{Top layer} (highest observed surface):
\begin{equation}
h^{\text{top}}_{i,j} \leftarrow (1 - \alpha_t)\, h^{\text{top}}_{i,j} + \alpha_t\, \max(\{z_k\}_{k \in \text{cell}})
\end{equation}

with $\alpha_t = 0.50$. The ground layer uses a low quantile to robustly estimate the walkable surface even when obstacle points fall in the same cell.

\subsection{Derived Terrain Quantities}

From the dual-layer heightmap we derive three quantities consumed by downstream modules:

\textbf{Surface normal.} Estimated via central differences on the ground layer:

\begin{equation}
\hat{\mathbf{n}}_{i,j} = \frac{1}{\|\mathbf{n}\|}\begin{bmatrix} -\partial h^{\text{ground}} / \partial x \\ -\partial h^{\text{ground}} / \partial y \\ 1 \end{bmatrix}, \quad \frac{\partial h^{\text{ground}}}{\partial x} \approx \frac{h^{\text{ground}}_{i,j+1} - h^{\text{ground}}_{i,j-1}}{2\delta}
\end{equation}

\textbf{Clearance.} Per-cell vertical gap between top and ground layers:

\begin{equation}
c_{i,j} = h^{\text{top}}_{i,j} - h^{\text{ground}}_{i,j}
\end{equation}

Cells with $c_{i,j} \geq c_{\text{lethal}} = 0.25$ m are classified as obstacles.

\textbf{Steppability (Foothold Quality).} A composite terrain quality metric $\sigma \in [0, 1]$ (0 = ideal, 1 = untraversable):

\begin{equation}
\sigma(\mathbf{p}) = 0.6\, \sigma_{\text{slope}}(\mathbf{p}) + 0.4\, \sigma_{\text{rough}}(\mathbf{p})
\end{equation}

where

\begin{equation}
\sigma_{\text{slope}} = \text{clip}\!\left(\frac{\|\nabla h^{\text{ground}}\|}{0.5},\; 0,\; 1\right), \quad \sigma_{\text{rough}} = \text{clip}\!\left(\frac{\text{std}_5(\mathbf{p})}{0.03},\; 0,\; 1\right)
\end{equation}

A slope magnitude of 0.5 ($\approx 27°$) saturates the slope term; surface roughness of 3 cm saturates the roughness term.

\subsection{Obstacle Cost Map}

Obstacle cells are inflated radially to create a smooth cost field for planning:

\begin{equation}
\mathcal{C}(\mathbf{p}) = \max_{(i,j) \in \mathcal{O}} \left[ \max\!\left(0,\; 1 - \frac{\|\mathbf{p} - \mathbf{p}_{ij}\|}{r_{\text{inflate}}}\right) \right]
\end{equation}

where $\mathcal{O}$ is the set of lethal cells and $r_{\text{inflate}} = 0.70$ m. The cost map applies temporal decay $\mathcal{C} \leftarrow 0.98\,\mathcal{C}$ each update cycle.

\section{GLOBAL PATH PLANNING}

\subsection{Terrain-Aware Weighted A*}

Given the robot's current position $\mathbf{p}_{\text{start}}$ and a goal $\mathbf{p}_{\text{goal}}$, we compute a global path on the cost map grid using weighted A*. The edge cost between adjacent cells incorporates both obstacle proximity and terrain traversability:

\begin{equation}
w(i \to j) = \|\mathbf{p}_i - \mathbf{p}_j\| \cdot \left(1 + w_{\text{obs}}\left(e^{4\,\mathcal{C}_j} - 1\right) + w_{\text{terr}}\,\sigma_j\right)
\end{equation}

where $w_{\text{obs}} = 100$, $w_{\text{terr}} = 5$, $\mathcal{C}_j$ is the obstacle cost, and $\sigma_j$ is the steppability score at cell $j$. Cells with $\mathcal{C} \geq 0.8$ are treated as impassable.

The heuristic is $h(\mathbf{p}) = 1.5\,\|\mathbf{p} - \mathbf{p}_{\text{goal}}\|$ (weighted A* with $\epsilon = 1.5$ for faster planning).

\subsection{Path Hysteresis}

To prevent left-right flip-flopping on successive replans, cells lying on the previous path receive a 50\% edge-cost discount:

\begin{equation}
w_{\text{hysteresis}}(i \to j) = \begin{cases} 0.5\, w(i \to j) & \text{if } j \in \mathcal{P}_{\text{prev}} \\ w(i \to j) & \text{otherwise} \end{cases}
\end{equation}

\subsection{Path Post-Processing}

The raw A* path is resampled to uniform arc-length spacing of 0.15 m, then smoothed via 40 iterations of Laplacian averaging with collision checking:

\begin{equation}
\mathbf{p}_k \leftarrow (1 - \alpha)\,\mathbf{p}_k + \frac{\alpha}{2}(\mathbf{p}_{k-1} + \mathbf{p}_{k+1}), \quad \alpha = 0.25
\end{equation}

Each smoothed candidate is accepted only if $\mathcal{C}(\mathbf{p}_k) < 0.8$.

\section{LOCAL TRAJECTORY OPTIMIZATION (MPPI)}

\subsection{Dynamics Model}

The MPPI local planner generates body-frame velocity commands $\mathbf{u} = [v_x, v_y, \omega_z]^T$ by optimizing over a finite horizon. We model the robot as a planar unicycle with first-order velocity dynamics:

\begin{align}
\dot{x} &= v_x \cos\psi - v_y \sin\psi, \quad \dot{y} = v_x \sin\psi + v_y \cos\psi, \quad \dot{\psi} = \omega_z \\
\dot{v}_x &= \frac{v_x^{\text{cmd}} - v_x}{\tau}, \quad \dot{v}_y = \frac{v_y^{\text{cmd}} - v_y}{\tau}, \quad \dot{\omega}_z = \frac{\omega_z^{\text{cmd}} - \omega_z}{\tau}
\end{align}

where $\tau = 0.4$ s is a velocity tracking time constant. The state is $\mathbf{s} = [x, y, \psi, v_x, v_y, \omega_z]^T \in \mathbb{R}^6$.

\subsection{Sampling and Optimization}

At each MPPI cycle, we maintain a nominal control sequence $\mathbf{U} = [\mathbf{u}_0, \ldots, \mathbf{u}_{H-1}] \in \mathbb{R}^{H \times 3}$ with horizon $H = 80$ steps. We draw $K = 400$ perturbation sequences with temporally-correlated AR(1) noise:

\begin{equation}
\boldsymbol{\epsilon}^{(k)}_t = \alpha\,\boldsymbol{\epsilon}^{(k)}_{t-1} + (1 - \alpha)\,\boldsymbol{\xi}^{(k)}_t, \quad \boldsymbol{\xi}^{(k)}_t \sim \mathcal{N}(\mathbf{0}, \boldsymbol{\Sigma})
\end{equation}

where $\alpha = 0.5$ and $\boldsymbol{\Sigma} = \text{diag}(\sigma_{v_x}^2, \sigma_{v_y}^2, \sigma_{\omega_z}^2)$. The nominal sequence is updated via importance-weighted averaging:

\begin{equation}
\mathbf{U} \leftarrow \sum_{k=1}^{K} w_k\, \mathbf{U}^{(k)}, \quad w_k = \frac{\exp\left(-\frac{1}{\lambda}(J_k - J_{\min})\right)}{\sum_{j=1}^{K}\exp\left(-\frac{1}{\lambda}(J_j - J_{\min})\right)}
\end{equation}

with temperature $\lambda = 8.0$. This is repeated for $I = 5$ iterations per cycle.

\subsection{Cost Function (Nav2-Style Critics)}

The total cost for rollout $k$ is a sum of eight critic terms:

\begin{equation}
J_k = J_{\text{path}} + J_{\text{heading}} + J_{\text{turn}} + J_{\text{progress}} + J_{\text{obs}} + J_{\text{step}} + J_{\text{slope}} + J_{\text{smooth}}
\end{equation}

\textbf{1) Path-Following Critic.} Penalizes cross-track error:

\begin{equation}
J_{\text{path}} = 80 \cdot \frac{1}{|\{t_s\}|}\sum_{t_s} \min_{\mathbf{p}_j \in \mathcal{P}} \|\mathbf{x}^{(k)}_{t_s} - \mathbf{p}_j\|
\end{equation}

\textbf{2) Path Angle Critic.} Aligns heading $\psi$ with path tangent $\hat{\mathbf{t}}$:

\begin{equation}
J_{\text{heading}} = 120 \cdot \frac{1}{|\{t_s\}|}\sum_{t_s} \left|\text{atan2}\!\left(\sin(\psi_{t_s} - \theta_{\hat{\mathbf{t}}}),\; \cos(\psi_{t_s} - \theta_{\hat{\mathbf{t}}})\right)\right|
\end{equation}

\textbf{3) Turn-in-Place Critic.} Prevents forward motion when heading error is large:

\begin{equation}
J_{\text{turn}} = 500 \cdot \frac{1}{|\{t_s\}|}\sum_{t_s} \left[\max(v_{x,t_s}, 0)\right]^2 \cdot \left[\max(|e_{\psi,t_s}| - 0.2, 0)\right]^2
\end{equation}

\textbf{4) Path Progress Critic.} Rewards reaching further along path arc length:

\begin{equation}
J_{\text{progress}} = 20 \cdot \frac{L_{\text{total}} - s_{\max}^{(k)}}{L_{\text{total}}}
\end{equation}

\textbf{5) Obstacle Critic.} Graduated penalty using cost map $\mathcal{C}$:

\begin{equation}
J_{\text{obs}} = \frac{1}{20}\sum_{t=0}^{H-1} P_t + 3 \cdot \max_t P_t, \quad P_t = \begin{cases} 300 & \text{if } \mathcal{C}_t \geq 1.0 \\ 150\,\mathcal{C}_t & \text{if } \mathcal{C}_t > 0.3 \\ 30\,\mathcal{C}_t & \text{otherwise} \end{cases}
\end{equation}

\textbf{6) Foothold-Quality-Aware (FQA) Steppability Critic.} Predicts foot locations and penalizes poor stepping surfaces:

\begin{equation}
\mathbf{p}_l^{\text{foot}}(t) = \begin{bmatrix} x_t \\ y_t \end{bmatrix} + \mathbf{R}(\psi_t)\,\mathbf{h}_l
\end{equation}

\begin{equation}
J_{\text{step}} = \frac{w_{\text{step}}}{4} \sum_{l=1}^{4} \frac{1}{H}\sum_{t=0}^{H-1} \sigma^2(\mathbf{p}_l^{\text{foot}}(t)), \quad w_{\text{step}} = 3.0
\end{equation}

\textbf{7) Slope Critic.} Penalizes steep terrain:

\begin{equation}
J_{\text{slope}} = 2.5 \cdot \frac{1}{H}\sum_{t=0}^{H-1} \|\nabla h^{\text{ground}}(\mathbf{x}_t)\|
\end{equation}

\textbf{8) Smoothness Critic.} Regularizes control changes:

\begin{equation}
J_{\text{smooth}} = 0.1 \cdot \frac{1}{H-1}\sum_{t=0}^{H-2} |\mathbf{u}_{t+1} - \mathbf{u}_t|
\end{equation}

\subsection{Adaptive Noise and Warm-Starting}

The noise standard deviation $\boldsymbol{\Sigma}$ is adapted based on context:

\begin{itemize}
\item \textbf{Near goal and on-track} ($d < 1.5$ m, cross-track $< 0.3$ m): $\boldsymbol{\sigma} = [0.08, 0.05, 0.40]$ (narrow, exploit)
\item \textbf{Stuck or near obstacle}: $\boldsymbol{\sigma} = [0.30, 0.10, 1.20]$ (wide yaw exploration)
\item \textbf{Default}: $\boldsymbol{\sigma} = [0.35, 0.05, 1.00]$
\end{itemize}

The first command $\mathbf{u}_0$ is acceleration-clamped before execution:

\begin{equation}
\mathbf{u}_0 \leftarrow \mathbf{u}_{\text{prev}} + \text{clip}(\mathbf{u}_0 - \mathbf{u}_{\text{prev}},\; -\mathbf{a}_{\max}\Delta t,\; \mathbf{a}_{\max}\Delta t)
\end{equation}

with $\mathbf{a}_{\max} = [2.0, 1.5, 10.0]$ in $[\text{m/s}^2, \text{m/s}^2, \text{rad/s}^2]$.

\section{CENTROIDAL MODEL PREDICTIVE CONTROL}

\subsection{Centroidal Dynamics}

We model the quadruped as a single rigid body with mass $m$ and rotational inertia $\mathbf{I}_{\text{com}}$ in the world frame. The centroidal state is

\begin{equation}
\mathbf{x} = \begin{bmatrix} \mathbf{p} \\ \boldsymbol{\Theta} \\ \dot{\mathbf{p}} \\ \boldsymbol{\omega} \end{bmatrix} \in \mathbb{R}^{12}
\end{equation}

where $\mathbf{p} \in \mathbb{R}^3$ is the center-of-mass position, $\boldsymbol{\Theta} = [\phi, \theta, \psi]^T$ are roll-pitch-yaw Euler angles, $\dot{\mathbf{p}} \in \mathbb{R}^3$ is COM velocity, and $\boldsymbol{\omega} \in \mathbb{R}^3$ is angular velocity.

The control input is the concatenated contact forces at four feet:

\begin{equation}
\mathbf{u} = \begin{bmatrix} \mathbf{f}_{\text{FL}} \\ \mathbf{f}_{\text{FR}} \\ \mathbf{f}_{\text{RL}} \\ \mathbf{f}_{\text{RR}} \end{bmatrix} \in \mathbb{R}^{12}
\end{equation}

The continuous-time dynamics are:

\begin{equation}
\dot{\mathbf{x}} = \mathbf{A}_c\,\mathbf{x} + \mathbf{B}_c(t)\,\mathbf{u} + \mathbf{g}_c
\end{equation}

with

\begin{equation}
\mathbf{A}_c = \begin{bmatrix} \mathbf{0} & \mathbf{0} & \mathbf{I}_3 & \mathbf{0} \\ \mathbf{0} & \mathbf{0} & \mathbf{0} & \mathbf{R}_z^T \\ \mathbf{0} & \mathbf{0} & \mathbf{0} & \mathbf{0} \\ \mathbf{0} & \mathbf{0} & \mathbf{0} & \mathbf{0} \end{bmatrix}, \quad \mathbf{B}_c(t) = \begin{bmatrix} \mathbf{0}_{3 \times 12} \\ \mathbf{0}_{3 \times 12} \\ \frac{1}{m}\begin{bmatrix}\mathbf{I}_3 & \mathbf{I}_3 & \mathbf{I}_3 & \mathbf{I}_3\end{bmatrix} \\ \mathbf{I}_{\text{com}}^{-1}\begin{bmatrix}[\mathbf{r}_1]_\times & [\mathbf{r}_2]_\times & [\mathbf{r}_3]_\times & [\mathbf{r}_4]_\times\end{bmatrix} \end{bmatrix}
\end{equation}

\begin{equation}
\mathbf{g}_c = \begin{bmatrix} \mathbf{0}_6 \\ [0, 0, -g]^T \\ \mathbf{0}_3 \end{bmatrix}
\end{equation}

\subsection{Discretization}

We discretize over a horizon of $N = 16$ steps with timestep $\Delta t = T_{\text{gait}}/N$. Using zero-order hold (ZOH):

\begin{equation}
\mathbf{A}_d, \mathbf{B}_{d,k} = \text{ZOH}(\mathbf{A}_c, \mathbf{B}_{c,k}, \Delta t)
\end{equation}

The discrete gravity term requires integration:

\begin{equation}
\mathbf{g}_d = \int_0^{\Delta t} e^{\mathbf{A}_c \tau}\,\mathbf{g}_c\,d\tau \approx \text{trapz}\!\left(\left\{e^{\mathbf{A}_c \tau_i}\,\mathbf{g}_c\right\}_{i=0}^{49}\right)
\end{equation}

\subsection{Convex QP Formulation}

The MPC solves:

\begin{equation}
\min_{\mathbf{x}_{1:N},\, \mathbf{u}_{0:N-1}} \sum_{k=0}^{N-1}\left[\|\mathbf{x}_k - \mathbf{x}_k^{\text{ref}}\|_{\mathbf{Q}}^2 + \|\mathbf{u}_k\|_{\mathbf{R}}^2\right]
\end{equation}

subject to:

\begin{align}
\mathbf{x}_{k+1} &= \mathbf{A}_d\,\mathbf{x}_k + \mathbf{B}_{d,k}\,\mathbf{u}_k + \mathbf{g}_d \\
\mathbf{f}_l &= \mathbf{0} \quad \text{(swing legs)} \\
(\hat{\mathbf{t}}_i - \mu\,\hat{\mathbf{n}}_l)^T \mathbf{f}_l &\leq 0, \quad i = 1, 2 \\
-\hat{\mathbf{n}}_l^T \mathbf{f}_l &\leq -f_{\min}
\end{align}

Cost weights: $\mathbf{Q} = \text{diag}(1, 1, 50, 10, 20, 1, 2, 2, 1, 1, 1, 1)$, $\mathbf{R} = 10^{-5}\mathbf{I}_{12}$.

\subsection{Terrain-Adaptive Friction Cones}

Surface normals $\hat{\mathbf{n}}_l$ from the heightmap are used to construct an orthonormal tangent basis via Gram-Schmidt:

\begin{align}
\hat{\mathbf{t}}_1 &= \frac{\hat{\mathbf{n}} \times \mathbf{a}}{\|\hat{\mathbf{n}} \times \mathbf{a}\|} \\
\hat{\mathbf{t}}_2 &= \frac{\hat{\mathbf{n}} \times \hat{\mathbf{t}}_1}{\|\hat{\mathbf{n}} \times \hat{\mathbf{t}}_1\|}
\end{align}

The friction pyramid faces are:

\begin{equation}
(\hat{\mathbf{t}}_i - \mu\,\hat{\mathbf{n}}_l)^T \mathbf{f}_l \leq 0, \quad (-\hat{\mathbf{t}}_i - \mu\,\hat{\mathbf{n}}_l)^T \mathbf{f}_l \leq 0, \quad i = 1, 2
\end{equation}

with $\mu = 0.8$ and $f_{\min} = 10$ N.

\section{REFERENCE TRAJECTORY GENERATION}

\subsection{Position and Velocity}

World-frame reference velocity is $\mathbf{v}_{\text{ref}} = \mathbf{R}_z(\psi)\,[v_x^{\text{cmd}}, v_y^{\text{cmd}}, 0]^T$. Position is extrapolated linearly with error clamping $\|\mathbf{p}^{\text{ref}} - \mathbf{p}_{\text{current}}\|_\infty \leq 0.1$ m.

\subsection{Terrain-Consistent Height}

\begin{equation}
z_k^{\text{ref}} = \max\!\left(h^{\text{ground}}(x_k, y_k),\; \bar{h}_{\text{hip},k}\right) + z_{\text{des}}
\end{equation}

\subsection{Slope-Aligned Orientation}

\begin{align}
\phi_k^{\text{ref}} &= -\kappa\,\arctan\!\left(\frac{n_y}{n_z}\right) \\
\theta_k^{\text{ref}} &= \kappa\,\arctan\!\left(\frac{n_x}{n_z}\right)
\end{align}

with blending factor $\kappa = 0.25$.

\section{GAIT SCHEDULING AND FOOTHOLD SELECTION}

\subsection{Trot Gait}

Contact phase for leg $l$ at time $t$:

\begin{equation}
c_l(t) = \begin{cases} 1 & \text{if } \left(\frac{t}{T} + \phi_l\right) \bmod 1 < D \\ 0 & \text{otherwise} \end{cases}
\end{equation}

with phase offsets $\boldsymbol{\phi} = [0, 0.25, 0.5, 0.75]$ and duty cycle $D = 0.8$.

\subsection{Touchdown Position}

\begin{equation}
\mathbf{p}_{\text{td}} = \mathbf{p}_{\text{hip}} + \mathbf{v}_{\text{des}}\,\tfrac{T_{\text{pred}}}{2} + k_p(\mathbf{p}_{\text{com}} - \mathbf{p}_{\text{des}}) + k_v(\dot{\mathbf{p}}_{\text{com}} - \mathbf{v}_{\text{des}}) + \boldsymbol{\delta}_{\text{yaw}} + \boldsymbol{\delta}_{\text{capture}}
\end{equation}

\textbf{Yaw rotation correction:}

\begin{equation}
\boldsymbol{\delta}_{\text{yaw}} = \begin{bmatrix} -\dot{\psi}\,T_{\text{pred}}\,r_y \\ \dot{\psi}\,T_{\text{pred}}\,r_x \\ 0 \end{bmatrix}
\end{equation}

\textbf{Roll-reactive capture point (LIPM):}

\begin{equation}
\omega_0 = \sqrt{\frac{g}{h_{\text{com}}}}, \quad x_{\text{cap}} = \frac{v_{\text{lat}}}{\omega_0} - \sin(\phi)\,h_{\text{com}}
\end{equation}

For right legs when falling rightward:

\begin{equation}
\boldsymbol{\delta}_{\text{capture}} = 0.6\,\mathbf{R}_z\,[0,\; x_{\text{cap}},\; 0]^T
\end{equation}

\subsection{Terrain-Aware Foothold Selection}

\textbf{Hard gates:} Slope $> 25°$, reachability $\notin [0.08, 0.35]$ m, height jump $> 0.10$ m, plane-fit residual $> 0.015$ m.

\textbf{Support polygon stability.} COM prediction with roll acceleration:

\begin{equation}
\mathbf{p}_{\text{com}}^{\text{pred}} = \mathbf{p}_{\text{com}} + (\dot{\mathbf{p}}_{\text{com}} + \dot{\mathbf{p}}_{\text{lat}})\,t_{\text{swing}} + \tfrac{1}{2}\,g\sin(\phi)\,t_{\text{swing}}^2\,\hat{\mathbf{y}}_{\text{world}}
\end{equation}

\textbf{Multi-objective scoring:}

\begin{equation}
S = w_d\,\bar{d} + w_s\,\bar{s} + w_r\,\bar{r} + w_f\,\bar{f} + w_g\,\bar{g} + w_{\text{stab}}(\phi)\,\bar{m}
\end{equation}

\textbf{Adaptive stability weight:}

\begin{equation}
w_{\text{stab}}(\phi) = 0.20 + 0.50\,\text{clip}\!\left(\frac{|\phi|}{10°} + \frac{|\dot{\phi}|}{1.0},\; 0,\; 1\right)
\end{equation}

\subsection{Swing Trajectory}

Minimum-jerk trajectory:

\begin{equation}
\mathbf{p}(s) = \mathbf{p}_0 + (\mathbf{p}_f - \mathbf{p}_0)\,\mu(s), \quad \mu(s) = 10s^3 - 15s^4 + 6s^5
\end{equation}

Height bump:

\begin{equation}
\Delta z(s) = h_{\text{sw}} \cdot 64\,s^3(1 - s)^3
\end{equation}

\section{LOW-LEVEL LEG CONTROLLER}

\subsection{Swing Phase}

Cartesian impedance with inverse-dynamics feedforward:

\begin{equation}
\mathbf{f}_{\text{swing}} = \mathbf{K}_p(\mathbf{p}_{\text{des}} - \mathbf{p}) + \mathbf{K}_d(\dot{\mathbf{p}}_{\text{des}} - \dot{\mathbf{p}}) + \boldsymbol{\Lambda}(\ddot{\mathbf{p}}_{\text{des}} - \dot{\mathbf{J}}\dot{\mathbf{q}})
\end{equation}

where $\boldsymbol{\Lambda} = (\mathbf{J}\mathbf{M}^{-1}\mathbf{J}^T)^{-1}$, $\mathbf{K}_p = 400\,\mathbf{I}_3$ N/m, $\mathbf{K}_d = 75\,\mathbf{I}_3$ Ns/m.

\subsection{Stance Phase}

Force tracking with foothold hold:

\begin{equation}
\mathbf{f}_{\text{stance}} = -\mathbf{f}_{\text{mpc}} + \mathbf{K}_p^{\text{st}}(\mathbf{p}_{\text{td}} - \mathbf{p}) + \mathbf{K}_d^{\text{st}}(-\dot{\mathbf{p}})
\end{equation}

with $\mathbf{K}_p^{\text{st}} = \text{diag}(300, 300, 0)$ N/m, $\mathbf{K}_d^{\text{st}} = \text{diag}(40, 40, 20)$ Ns/m.

\subsection{Joint Torques}

\begin{equation}
\boldsymbol{\tau} = \mathbf{J}_{\text{leg}}^T\,\mathbf{f} + (\mathbf{C}\dot{\mathbf{q}} + \mathbf{g})_{\text{leg}}
\end{equation}

Joint torques are saturated at 90\% of hardware limits.

\section{INTEGRATION: TERRAIN DATA FLOW}

\begin{center}
\begin{tabular}{|l|l|l|}
\hline
\textbf{Consumer} & \textbf{Terrain Query} & \textbf{Effect} \\
\hline
A* planner & Steppability $\sigma$ per cell & Routes around poor surfaces \\
MPPI planner & Steppability at foothold predictions & Penalizes rough/steep terrain \\
MPPI planner & Obstacle cost map $\mathcal{C}$ & Avoids collisions \\
MPPI planner & Terrain gradient $\nabla h$ & Penalizes steep slopes \\
COM trajectory & Ground height $h^{\text{ground}}$, normal $\hat{\mathbf{n}}$ & Terrain-consistent COM height/orientation \\
Centroidal MPC & Surface normal $\hat{\mathbf{n}}_l$ per foot & Rotates friction cone constraints \\
Foothold selector & Slope, roughness, residual, height & Multi-criteria terrain quality \\
Swing trajectory & Touchdown height $h^{\text{ground}}$ & Sets correct foot landing height \\
\hline
\end{tabular}
\end{center}

\end{document}
